\chapter*{Introduction to Personal Financial Planning}
\addcontentsline{toc}{chapter}{Introduction to Personal Financial Planning}

Second year of university is a strange time to be thinking about retirement. But that is roughly where this report starts---with the question of what someone at 20 needs to do today to have a comfortable life at 60, and what the 40 years in between should look like. The gap between now and then is large enough that it is easy to ignore. It is also exactly long enough that the decisions made in it compound into something either very good or very hard to undo.

This report is a personal financial plan for me, Vũ Đức Huy, a second-year Banking and Finance student at Foreign Trade University, Hanoi. I live in a family-owned 72~m\textsuperscript{2} apartment where rent, electricity, and water are covered by my family, which keeps my cost of living much lower than most students in Hanoi. My income runs from three sources at the same time: a 3,000,000~VND internship allowance from kNowHow Consultant, average freelance software income of 2,500,000~VND, and a 4,000,000~VND monthly family allowance---9,500,000~VND total. Two people shape how I think about discipline and preparation: Max Verstappen (Formula~1) and Chou Tien Chen (Badminton). Not for their results, but for how they work.

The report covers ten phases: spending analysis, goal-setting, personal financial statements, income tax, budgeting, credit strategy, vehicle purchase, stock investment, insurance, and retirement planning. All numbers come from a real 30-day transaction log (May~7 to June~4, 2026) exported from my ABBank app. The path traced runs from 9,500,000 VND per month now to a projected 300,000,000 VND per month as a BCG Partner or technology founder, with a retirement target at age~60.
